% !TEX root = main.tex

% Font options: use `times` or `arial` with the class.
\documentclass[conference,times]{unisatex}
\unisatexoverridecommandlockouts

\usepackage{tikz-cd}
\usepackage{multicol}

\title{Título do Artigo\unisatexsubtitle{Modelo institucional para artigos}}

\author{
    \unisatexauthorblockN{1\textsuperscript{st} Given Name Surname}
    \unisatexauthorblockA{
        \textit{Department of Computer Engineering} \\
        \textit{UniSATC} \\
        Criciúma, Brasil \\
        email address or ORCID
    }

    \and

    \unisatexauthorblockN{2\textsuperscript{nd} Given Name Surname}
    \unisatexauthorblockA{
        \textit{Department of Computer Engineering} \\
        \textit{UniSATC} \\
        Criciúma, Brasil \\
        email address or ORCID
    }
}

\begin{document}
\maketitle

\begin{abstract}
Este documento apresenta uma base inicial para artigos acadêmicos da UniSATC em formato inspirado no modelo de conferências do IEEE~\cite{ieeetran}. A classe local preserva o arquivo \texttt{IEEEtran.cls} como dependência e concentra as personalizações institucionais em \texttt{unisatex.cls}.
\end{abstract}

\begin{keywords}
UniSATC, modelo LaTeX, conferência, artigo acadêmico, IEEEtran.
\end{keywords}

\section{Introdução}
Este arquivo demonstra o uso do modelo \texttt{unisatex.cls}. A estrutura segue a organização típica de artigos de conferência: título, autores, resumo, palavras-chave, seções numeradas, figuras, tabelas, equações e referências.

\section{Preparação do Artigo}
Escreva o conteúdo principal antes de ajustar detalhes de formatação. O objetivo da classe é manter a consistência visual e evitar mudanças manuais nas margens, nas colunas e nos espaçamentos.

\subsection{Equações}
As equações podem ser numeradas normalmente:
\begin{equation}
E = mc^2.
\end{equation}

\subsection{Figuras e Tabelas}
Insira figuras e tabelas após serem citadas no texto. Tabelas devem ter legenda acima e figuras devem ter legenda abaixo, seguindo a prática comum do IEEE~\cite{ieeeconference}.

\section{Conclusão}
O modelo pode ser expandido com regras específicas da UniSATC, como logotipo, identificação de evento, rodapé institucional, normas de submissão e exemplos bibliográficos.

\bibliographystyle{IEEEtran}
\bibliography{references}

\end{document}